\section{Regime-map definitions and study design}

\subsection{Exactness and support}

Let $x\in\mathcal{X}$ be an instance, let $\theta$ specify objective weights, and let $\mathcal{C}(x)$ be the feasible compiler configurations.  The exact cost is
\begin{equation}
 C^\star(x;\theta)=\min_{c\in\mathcal{C}(x)} C(c;x,\theta).
\end{equation}
A deterministic reference compiler returns cost $C_D(x;\theta)$.  Its exact region is
\begin{equation}
 \mathcal{R}_D(\theta)=\{x:C_D(x;\theta)=C^\star(x;\theta)\}.
\end{equation}
Equality is always indexed by the model, domain and objective.  A finite census, an all-size theorem and a frozen panel are reported as different evidence classes.

For models with a support coordinate $s(c)$, define the intrinsic support number by
\begin{equation}
 \kappa(\theta)=\min\{k:\forall x\;\exists c^\star\text{ with }s(c^\star)\leq k\}.
\end{equation}
A proof of sufficiency gives an upper bound on $\kappa$.  Calling the bound intrinsic also requires tightness, such as infeasibility below it.

\subsection{Constructive mechanisms and proof-validity regions}

When the reference is inexact, a mechanism consists of a feasible witness, a strict cost reduction, and the structural edit that produces it.  We use mechanism names only after these three parts are checked.  A growing mechanism vocabulary does not weaken a separate theorem about the full feasible family.

A normalization proof may depend on $\theta$.  Its \emph{proof-validity region} is the set of objective weights for which every local replacement is non-increasing and the analytic composition remains valid.  Outside that set, the certificate is silent.  A counterexample or a necessity theorem must be established separately.

\subsection{Feature determination}

Let $\phi(x)$ be a frozen feature vector and $y(x)=1[x\in\mathcal{R}_D]$.  A feature cell is the set of instances sharing one vector.  A cell is mixed if it contains both labels.  The minimum error attainable by any deterministic classifier using only $\phi$ on a finite domain is
\begin{equation}
 E_{\min}(\phi)=\sum_v \min\{n_0(v),n_1(v)\},
\end{equation}
where $n_b(v)$ counts label $b$ in cell $v$.  Thus, $E_{\min}>0$ is a representation-level impossibility result.  Conversely, $E_{\min}=0$ only proves determination on the registered domain.  Cell count and singleton count reveal whether that determination is compressive or nearly memorizing.

\subsection{Typed regime-map schema}

For a declared model $M$, objective $\theta$ and domain $\mathcal X$, we record
\begin{equation}
 \mathfrak M(M,\theta,\mathcal X)=
 (\mathcal R_D,\mathcal W,\kappa,B,\Omega,\phi,E_{\min},\mathcal T;\alpha).
 \label{eq:typed-map}
\end{equation}
Here $\mathcal W$ is the set of verified constructive gap witnesses, $B$ is a proof-derived support ceiling, $\Omega$ is the objective region in which its proof is valid, $\phi$ is a fixed feature representation, $\mathcal T$ is a prospective-transfer record, and $\alpha$ states the evidence class (all-size proof, complete finite census, frozen panel or bounded search).  Each coordinate carries one of three statuses: \emph{value}, \emph{unknown} (defined but not established) or \emph{not applicable} (the coordinate is absent from the model).  An unknown value is never replaced by another coordinate; in particular, $B$ does not become $\kappa$ without tightness, and $E_{\min}=0$ does not populate $\mathcal T$.  Two numerical coordinates are comparable only when their mathematical type, model family, objective and domain agree.  Across different models we compare logical role and evidence class, not numerical magnitude.

\subsection{Model register}

We represent an $n$-qubit Pauli word, modulo global phase, by its binary symplectic vector.  Its Hamming support is $\operatorname{wt}(P)$ and $\langle P,Q\rangle\in\mathbb F_2$ denotes commutation parity.  Full feasible sets and objectives are stated in Supplementary Section~S1.

The shared-tag Pauli model contains three blocks, each with two nonidentity frames $R_{j0},R_{j1}$ satisfying $\langle R_{j0},R_{j1}\rangle=1$.  A shared Pauli tag $S$ gives the same two syndrome bits in every block, and those bits must differ.  The objective charges frame, shared-tag and restored-target support, with one central frame per block weighted two and the other weighted four.  The support coordinate is $\max_{j,k}\operatorname{wt}(R_{jk})$.

The rank-two dependent-triple model contains two blocks.  In block $b$, two anticommuting independent frames $R_{b0},R_{b1}$ determine $R_{b2}=R_{b0}R_{b1}$.  Shared tags $S_0,S_1$ assign the same distinct nonzero two-bit labels to the two independent frames in both blocks.  The objective weights central frame, noncentral frame, tag and restoration support by $(t_c,t_{nc},t_{\mathrm{tag}},t_r)$.  The support coordinate is the maximum support of an independent generator.

The six-term linear-combination model takes six nonidentity Pauli words, with repetition allowed, and compares the six-singleton reference encoding with all 203 set partitions.  Each block chooses whether to extract its common nonidentity Pauli factor, and index ancillas may be shared or dedicated.  The frozen objective is the sum of selection support, preparation-tree nodes and register width.  Coefficients and amplitude-dependent preparation costs are outside the model.

The state-preparation model uses $H$, $S$, $S^\dagger$ and controlled-NOT gates with costs $1,1,1,3$.  The reference is a deterministic Gaussian-elimination-style stabilizer synthesis.  Exact costs are shortest-path distances from $|0\cdots0\rangle$ in the complete stabilizer-state graph.  Complete graph sizes are 6, 60, 1,080 and 36,720 for $n=1,2,3,4$, respectively.  Complete analyses use $n\leq3$; the $n=4$ study uses a frozen 120-state panel.

\subsection{Verification and prospective discipline}

All-size results combine written compositions with complete finite checks of the local obligations.  Complete finite-domain results enumerate every admitted state or instance.  Selected calculations were rerun from a clean environment, and separate verification code checked the stated theorem fields, finite-domain arithmetic and manuscript values without invoking the scientific calculations.  These are same-team consistency checks, not external replication.

For the state-preparation transfer test, panel seed 20260821 generated 120 unique four-qubit states by applying 24 frozen random gates to $|0000\rangle$.  Exact labels came from Dijkstra search over all 36,720 stabilizer states.  The later representation concatenated 33 schedule/tensor coordinates, 53 reference-synthesis path coordinates and 41 sign-aware state coordinates.  Index parity defined two folds.  For each fold, a feature cell in the other fold predicted its majority label (ties negative); an unseen cell predicted negative.  Coverage means exact feature-vector equality with a training-fold cell.  The feature construction and criterion were fixed before this evaluation.  The shuffle null permuted the 120 labels 200 times with Python seed 20260828, preserved features and folds, and reran the identical lookup.  Its empirical probability is the fraction of permutations with error count no larger than the observed count.  Supplementary Sections~S5--S6 give the complete coordinate definitions, deterministic reference compiler, freeze order and null algorithm.
